\section{Introducción}

En las organizaciones que desarrollan soluciones tecnológicas para la industria, la documentación representa una parte importante del conocimiento operativo. Manuales, instructivos, procedimientos, políticas, reportes y especificaciones contienen información que debe consultarse para realizar actividades de manera uniforme. Sin embargo, la existencia de los documentos no garantiza que la información pueda localizarse con rapidez. Cuando los archivos se encuentran distribuidos en diferentes carpetas, tienen formatos distintos o utilizan una terminología que no coincide con la pregunta del usuario, la consulta puede convertirse en una actividad lenta y dependiente de la experiencia de una persona.

En un entorno de innovación industrial, esta situación adquiere una importancia adicional. Las soluciones se diseñan, prueban y modifican con frecuencia, por lo que la información técnica puede cambiar durante el ciclo de vida de un proyecto. Un integrante que necesita conocer un procedimiento no solo requiere una respuesta rápida; también necesita saber de qué documento provino la información, qué página o sección la contiene y si el material consultado corresponde a la versión vigente. La trazabilidad permite revisar la respuesta, corregir errores y evitar que una recomendación generada automáticamente se interprete como una instrucción oficial sin respaldo.

La inteligencia artificial (IA) generativa ofrece una forma natural de interactuar con la información mediante preguntas en lenguaje cotidiano. Un modelo de lenguaje puede interpretar la intención del usuario, resumir un texto y organizar una explicación. No obstante, si se consulta el modelo sin información específica de la organización, puede responder con conocimiento general, mezclar conceptos o producir una afirmación que parece correcta pero no está sustentada por los procedimientos internos. En consecuencia, la aplicación de inteligencia artificial a la documentación técnica requiere una arquitectura que combine generación, recuperación de información y control de fuentes \citep{lewis2020rag,gao2023rag}.

El presente proyecto propone la base de un copiloto agéntico basado en recuperación aumentada por generación (RAG, por sus siglas en inglés) para documentación técnica y procedimientos de planta. La solución integra una interfaz web, una interfaz de programación de aplicaciones (API, por sus siglas en inglés) propia y un flujo de inteligencia artificial. La interfaz permite que una persona escriba una pregunta y observe la conversación. El backend recibe la solicitud, valida sus datos y se comunica con el proveedor externo de IA. La arquitectura también contempla un recuperador semántico que buscará fragmentos relevantes en una colección documental, construirá el contexto de la consulta y devolverá una respuesta acompañada de fuentes. Este enfoque se fundamenta en la combinación de recuperación y generación descrita en la literatura especializada \citep{lewis2020rag,gao2023rag}.

En esta etapa se desarrolla la base necesaria para que los componentes puedan integrarse sin depender de una sola tecnología de interfaz o de un único proveedor de modelos. El frontend se implementa con HTML (HyperText Markup Language, lenguaje de marcado de hipertexto), CSS (Cascading Style Sheets, hojas de estilo en cascada) y JavaScript, mientras que el backend se construye con Python y FastAPI. La comunicación se realiza mediante HTTP (HyperText Transfer Protocol, protocolo de transferencia de hipertexto) y JSON (JavaScript Object Notation, notación de objetos de JavaScript). FastAPI proporciona validación de solicitudes y documentación interactiva de la API mediante OpenAPI \citep{fastapiDocs}. El proveedor de IA se consume desde el servidor para evitar que la clave de acceso quede expuesta en el navegador. La recuperación documental se define como una responsabilidad separada, de manera que los documentos, embeddings y metadatos puedan conectarse posteriormente al agente.

El proyecto se desarrolla en colaboración con el Centro de Innovación Industrial en Inteligencia Artificial, CII.IA, ubicado en Monterrey, Nuevo León, México. La organización comunica una orientación hacia la aplicación y el escalamiento de soluciones de inteligencia artificial para la industria, contexto que sustenta la pertinencia de una herramienta de consulta documental \citep{ciiia2026,ciiia2026b}. El resultado esperado no es únicamente una pantalla de conversación, sino una base técnica y documental para una herramienta que pueda crecer con el corpus institucional. A partir de esta base será posible incorporar criterios de similitud, reglas para detectar falta de contexto, fuentes consultables y pruebas que comparen las respuestas del agente con los documentos originales.

\subsection{Antecedentes}

La consulta de información ha pasado por diferentes etapas. En un primer momento, el usuario debía conocer la ubicación física o digital del documento y revisarlo manualmente. Posteriormente, los sistemas de archivos y las herramientas de búsqueda permitieron localizar palabras dentro de varios documentos. Aunque estos mecanismos reducen el tiempo de consulta, dependen de que el usuario utilice exactamente las mismas palabras que aparecen en el material. Una persona puede preguntar ``¿qué se debe hacer antes de iniciar el equipo?'' aunque el documento utilice una expresión como ``verificación previa a la puesta en operación''.

La búsqueda semántica intenta resolver esta diferencia al comparar el significado aproximado de la pregunta con el significado de los fragmentos documentales. Para ello, cada fragmento se transforma en un embedding y se almacena con información que permite rastrear su origen. El modelo generativo interviene después de la recuperación y utiliza el contexto seleccionado para elaborar una respuesta comprensible. Esta combinación se conoce como recuperación aumentada por generación.

El crecimiento de los modelos de lenguaje ha impulsado interfaces conversacionales para consultar bases de conocimiento. Sin embargo, un chatbot que no tenga conexión con el corpus de la organización no es suficiente para una consulta técnica controlada. Puede utilizar información que no está vigente, responder de forma general o no distinguir entre una recomendación y un procedimiento autorizado. La propuesta de este proyecto parte de esta diferencia: la conversación es el medio de interacción, pero los documentos deben ser la base de la respuesta.

\subsubsection{Antecedentes tecnológicos}

La arquitectura se apoya en tecnologías web y de servicios. HTML organiza los elementos de la pantalla, CSS define la distribución visual y JavaScript controla eventos, solicitudes y actualización del contenido. FastAPI permite declarar rutas, modelos de datos y validaciones para el servicio. Uvicorn ejecuta la aplicación ASGI (Asynchronous Server Gateway Interface, interfaz de puerta de enlace de servidor asíncrono) durante el desarrollo. La comunicación entre ambos extremos se realiza con solicitudes HTTP y objetos JSON.

Para la parte de inteligencia artificial se considera el consumo de un modelo externo mediante HTTPS (HyperText Transfer Protocol Secure, protocolo seguro de transferencia de hipertexto). El agente puede utilizar componentes de LangChain para preparar prompts y conectar un modelo con un recuperador, y LangGraph puede representar las decisiones del flujo mediante estados. PostgreSQL con capacidades vectoriales se contempla para almacenar embeddings y metadatos. Estas tecnologías se integran por capas, evitando que la lógica de recuperación quede mezclada con la presentación de la interfaz.

La seguridad básica también forma parte de los antecedentes del proyecto. Las claves de API deben existir únicamente en el entorno del backend, normalmente mediante un archivo `.env` que no se publica. El servidor debe validar la entrada antes de enviarla al proveedor y el frontend no debe conocer las credenciales. CORS (Cross-Origin Resource Sharing, intercambio de recursos de origen cruzado) se utiliza para controlar la comunicación entre el origen del frontend y el origen del backend durante el desarrollo.

\subsubsection{Antecedentes de la empresa}

El Centro de Innovación Industrial en Inteligencia Artificial, CII.IA, es una organización ubicada en Monterrey, Nuevo León, México, orientada a la aplicación de inteligencia artificial en contextos empresariales e industriales. Su trabajo relaciona la identificación de necesidades con el diseño, prototipado, validación y posible escalamiento de soluciones tecnológicas. Este contexto requiere que el conocimiento generado durante los proyectos pueda organizarse y consultarse de forma eficiente.

La creación de un copiloto documental se relaciona con esa necesidad porque convierte procedimientos y documentos técnicos en una fuente consultable mediante lenguaje natural. La herramienta no pretende reemplazar la responsabilidad de quien toma decisiones ni sustituir la validación de un experto. Su función inicial es apoyar la localización, comprensión y trazabilidad de la información, manteniendo visible el origen de los fragmentos utilizados.

El proyecto también es congruente con un ambiente de innovación porque utiliza una arquitectura que puede evolucionar. En una primera versión se valida la comunicación entre frontend, backend y modelo. En etapas posteriores se pueden incorporar nuevos documentos, actualizar embeddings, agregar reglas de evaluación y adaptar el flujo a otras áreas o tipos de procedimiento.

\subsection{Planteamiento del problema}

El problema se presenta cuando una persona necesita consultar una indicación técnica y no tiene un mecanismo uniforme para localizarla. La información puede estar dispersa en documentos distintos, contener versiones previas o emplear términos diferentes a los de la pregunta. La búsqueda manual consume tiempo y puede producir interpretaciones diferentes entre usuarios. Una búsqueda literal no siempre identifica fragmentos que expresan la misma idea con otras palabras.

La incorporación directa de un modelo de lenguaje tampoco resuelve por sí sola el problema. Aunque el modelo puede redactar una respuesta clara, no existe garantía de que utilice la documentación interna, de que conozca la versión vigente o de que pueda señalar la página donde obtuvo la información. En un procedimiento de planta, esta falta de trazabilidad representa una limitación importante.

Por lo tanto, el problema de investigación y desarrollo se formula de la siguiente manera: ¿cómo diseñar e implementar la base de un copiloto web que facilite la consulta de documentación técnica mediante inteligencia artificial, se comunique de manera segura con un backend, utilice un flujo preparado para recuperación semántica y devuelva respuestas con fuentes verificables?

El problema puede dividirse en cuatro aspectos. El primero es la interacción: se necesita una interfaz sencilla para escribir preguntas y leer respuestas. El segundo es la integración: se necesita una API que valide y coordine la solicitud. El tercero es el conocimiento: se necesita un corpus con fragmentos, embeddings y metadatos. El cuarto es el control: se necesitan reglas para indicar falta de información, manejar errores y evitar que una respuesta generada se presente como un procedimiento oficial sin evidencia.

\subsection{Objetivo general}

Desarrollar la base funcional de un copiloto agéntico RAG para consultar documentación técnica y procedimientos de planta, integrando una interfaz web, una API con FastAPI y un modelo externo de inteligencia artificial, con una arquitectura preparada para recuperar fragmentos documentales, conservar sus metadatos y presentar fuentes verificables.

\subsection{Objetivos específicos}

\begin{enumerate}
    \item Analizar la necesidad de consulta documental y delimitar el alcance funcional de la primera versión.
    \item Identificar los actores, entradas, salidas y restricciones que intervienen en una consulta al copiloto.
    \item Diseñar una interfaz web que permita escribir preguntas, enviar solicitudes y visualizar respuestas.
    \item Implementar una API con FastAPI para recibir, validar y responder solicitudes de conversación.
    \item Configurar CORS y el intercambio de datos en formato JSON entre el frontend y el backend.
    \item Integrar un proveedor externo de inteligencia artificial mediante variables de entorno protegidas.
    \item Definir el flujo de validación, recuperación, construcción de contexto y generación de respuestas.
    \item Establecer un contrato para fragmentos documentales que incluya identificador, contenido, archivo, página, sección y puntuación.
    \item Aplicar reglas para que el agente indique cuando no cuenta con información suficiente y evite inventar procedimientos.
    \item Preparar pruebas de comunicación entre frontend, backend, recuperador y modelo externo.
    \item Documentar la arquitectura, las decisiones técnicas, las limitaciones y las actividades pendientes.
\end{enumerate}

\subsection{Justificación}

Un copiloto para documentación técnica puede reducir el tiempo que se dedica a localizar información y facilitar el acceso al conocimiento organizacional. En lugar de depender de que el usuario recuerde el nombre de un archivo o revise manualmente varias carpetas, la interfaz permite formular una pregunta de manera natural. El resultado puede resumir el contenido y dirigir al usuario hacia el documento que debe revisar.

La propuesta es especialmente útil cuando la documentación contiene procedimientos repetitivos, definiciones técnicas o instrucciones que deben consultarse con frecuencia. La recuperación semántica puede identificar fragmentos relacionados aunque la redacción de la pregunta no coincida literalmente con la redacción del documento. La inclusión de metadatos permite que la respuesta sea revisable y que un responsable confirme la información antes de utilizarla en una actividad crítica.

Desde el punto de vista de desarrollo, el proyecto permite integrar conocimientos de frontend, backend, APIs, bases de datos, procesamiento de lenguaje natural y evaluación de inteligencia artificial. La separación de componentes facilita que el diseño visual pueda cambiar sin alterar el servicio, que el proveedor de IA pueda sustituirse y que la fuente documental pueda actualizarse sin reconstruir toda la aplicación.

La propuesta también es justificable por su enfoque de control. No se plantea un chatbot abierto que responda cualquier tema, sino un copiloto limitado a un dominio documental. Cuando el contexto no sea suficiente, el sistema debe comunicar esa condición. Esta decisión es importante porque una respuesta negativa y verificable es preferible a una respuesta extensa sin respaldo.

\subsection{Alcances y limitaciones}

Para delimitar con claridad la primera versión del proyecto, los alcances y las limitaciones se presentan de forma independiente.

\subsubsection{Alcances}

\begin{itemize}
    \item Analizar la necesidad de consulta de documentación técnica y definir el flujo general del copiloto.
    \item Mantener una interfaz web para que el usuario escriba preguntas, envíe solicitudes y consulte las respuestas recibidas.
    \item Implementar una API con FastAPI para validar la solicitud de conversación y devolver una respuesta en formato JSON.
    \item Configurar la comunicación entre frontend y backend mediante CORS, conservando la clave del proveedor de IA únicamente en el servidor.
    \item Integrar un modelo externo de inteligencia artificial mediante variables de entorno y un contrato de respuesta controlado.
    \item Definir la arquitectura RAG, incluyendo la consulta, la recuperación de contexto, la construcción del prompt y la generación de la respuesta.
    \item Especificar los metadatos mínimos de cada fragmento: \texttt{chunk\_id}, contenido, archivo de origen, página, sección y puntuación.
    \item Preparar la integración con la colección de documentos técnicos, sus embeddings, el almacenamiento vectorial y la búsqueda semántica entregada por el Integrante 1.
    \item Incluir fuentes y estados de respuesta para que el resultado pueda revisarse y distinguirse de una respuesta sin evidencia.
    \item Documentar la tabla comparativa, las decisiones de diseño, las pruebas de comunicación y las actividades pendientes.
\end{itemize}

\subsubsection{Limitaciones}

\begin{itemize}
    \item La respuesta depende de la disponibilidad, cuota, autenticación y modelo seleccionado en el proveedor externo de inteligencia artificial.
    \item La calidad de la respuesta depende de que los documentos estén completos, vigentes, correctamente fragmentados y acompañados de sus metadatos.
    \item La búsqueda semántica no garantiza por sí sola que un fragmento contenga toda la respuesta; por ello se requiere un umbral de similitud y validación del contexto.
    \item La colección documental, la carga de embeddings y el endpoint de búsqueda deben validarse con los datos autorizados y entregados por el Integrante 1.
    \item El sistema no sustituye la revisión de una persona responsable cuando la consulta se relaciona con seguridad, calidad, operación o una decisión crítica.
    \item La primera versión se limita al dominio de documentación técnica y procedimientos de planta; no se considera un chatbot de propósito general.
    \item El cambio de proveedor, la actualización del corpus y el despliegue en producción requieren configuraciones y pruebas adicionales.
\end{itemize}

\subsection{Organización del documento de tesina}

El Capítulo 1 presenta la introducción, los antecedentes generales, los antecedentes tecnológicos y de la empresa, el planteamiento del problema, los objetivos, la justificación, los alcances y las limitaciones. El Capítulo 2 desarrolla los antecedentes teóricos, los conceptos de inteligencia artificial generativa, RAG, embeddings, agentes, APIs, bases vectoriales y evaluación, además de los temas tentativos del marco teórico.

El Capítulo 3 presenta el estado del arte y los trabajos relacionados con la consulta de información, además de la tabla comparativa de alternativas. El Capítulo 4 conserva la sección destinada al desarrollo del proyecto. El Capítulo 5 se reserva para describir el sistema propuesto y su arquitectura. El Capítulo 6 contiene la estructura para documentar la implementación y las pruebas. El Capítulo 7 conserva el espacio para conclusiones y trabajo futuro. Las secciones de agradecimientos, resumen y summary se dejan vacías en esta versión para completarlas posteriormente.
